\documentclass[12pt,a4paper]{article}

\usepackage[utf8]{inputenc}
\usepackage[T1]{fontenc}
\usepackage[english]{babel}
\usepackage{setspace}
\usepackage{csquotes}
\usepackage{geometry}
\usepackage{array}
\usepackage{booktabs}
\usepackage{longtable}
\usepackage{tabularx}
\usepackage{xcolor}
\usepackage{listings}
\usepackage{hyperref}
\usepackage{graphicx}
\usepackage{fancyhdr}
\usepackage{enumitem}
\usepackage{textcomp}
\usepackage{xurl}

\setlength{\emergencystretch}{3em}
\urlstyle{same}

\geometry{left=3cm,right=2.5cm,top=2.5cm,bottom=2.5cm}
\hypersetup{
  colorlinks=true,
  linkcolor=black,
  urlcolor=blue,
  citecolor=black,
  pdftitle={ESP32 PV Surplus Load Controller Documentation},
  pdfauthor={Enrico Brodtmann}
}

\pagestyle{fancy}
\fancyhf{}
\lhead{ESP32 PV Surplus Load Controller}
\rhead{\thepage}
\setlength{\headheight}{14pt}
\renewcommand{\arraystretch}{1.2}
\setlist[itemize]{topsep=4pt,itemsep=2pt}
\setlist[enumerate]{topsep=4pt,itemsep=2pt}

\lstset{
  basicstyle=\ttfamily\small,
  breaklines=true,
  frame=single,
  columns=fullflexible,
  keepspaces=true,
  showstringspaces=false
}

\newcommand{\code}[1]{\texttt{\detokenize{#1}}}
\newcommand{\statusdone}{implemented}
\newcommand{\statuspartial}{partly implemented}
\newcommand{\statusfuture}{planned}
\newcommand{\statusoutofscope}{out of scope}

\begin{document}

\begin{titlepage}
    \centering
    \vspace*{1cm}

    \IfFileExists{hsb-logo_standardvariante.jpg}{%
        \includegraphics[width=0.45\textwidth]{hsb-logo_standardvariante.jpg}\\[1.5cm]
    }{%
        \vspace*{2.2cm}
    }

    {\Large Hochschule Bremen}\\[0.3cm]
    {\large Faculty 4 -- Electrical Engineering and Computer Science}\\[2cm]

    {\huge \textbf{ESP32 PV Surplus Load Controller}}\\[0.5cm]
    {\LARGE Documentation of a Local-First\\[0.3cm]
    PV Surplus Load Control Prototype\\[0.3cm]
    with Tactile UI and Home Assistant Integration}\\[2.5cm]

    \vfill
    \begin{flushleft}
        \textbf{Name / Student ID:}\\[0.2cm]
        Enrico Brodtmann (5265217)\\[0.5cm]
        \textbf{Degree Program:} Media Informatics B.Sc.\\
        \textbf{Module:} Microcontrollers and Applications\\
        \textbf{Lecturer:} C. Flores\\
        \textbf{Task:} Project documentation and prototype verification\\
        \textbf{Date:} June 29, 2026\\
        \textbf{Target specification:} \code{Microcontroler_Esp32_englisch_final.pdf}\\
        \textbf{Firmware:} v1 -- PlatformIO / Arduino ESP32\\
    \end{flushleft}
\end{titlepage}

\newpage
\tableofcontents
\newpage

\section{Abstract}

This project implements a proof-of-concept controller that uses an ESP32 to switch
a high-power load when photovoltaic surplus power is available. The ESP32 does not
switch mains voltage directly. Instead, it drives a low-voltage switching path which
can energize a solid-state relay and, through that relay, the coil of a contactor.
The contactor is the component intended to switch the real load.

The important design goal is that the actual control decision is made locally on
the ESP32. Home Assistant, MQTT and WiFi are used as useful interfaces for
monitoring, configuration and for providing the current surplus value in the
default build. They are not allowed to directly energize the output and cannot
override a local fault or local OFF state.

This documentation follows the exam requirement described in
\code{b_progress_and_exams.pdf}: the documentation should show how close the
implemented project came to its own target specification. Therefore, the document
does not only describe the program, but also compares the current implementation
against the original concept, explains the most important design decisions, and
defines verification steps for the prototype.

\newpage
\section{Project Goal and Scope}

\subsection{Original goal}

The target specification defines an ESP32-based control unit for using temporary
PV surplus energy. Typical loads are heating elements or other high-power
resistive consumers. These loads exceed the direct switching capability of a
microcontroller, so the ESP32 is only responsible for the safe low-voltage control
signal.

The core requirements from the target specification are:

\begin{itemize}
  \item receive or measure a value that represents available PV surplus power,
  \item make the switching decision locally on the ESP32,
  \item provide a switching output for an external relay, SSR or contactor,
  \item support the modes \code{OFF}, \code{MANUAL} and \code{AUTO},
  \item provide a real tactile local user interface,
  \item indicate status locally with LEDs and optionally a display,
  \item transmit status information through WiFi to Home Assistant or Node-RED,
  \item allow manual override at the device.
\end{itemize}

\subsection{Implemented prototype scope}

The implemented firmware focuses on the basic prototype: local control logic,
fault handling, a low-voltage switching output, tactile operation, temperature
monitoring and Home Assistant integration. The final expansion stage from the
concept document, namely integrated three-phase measurement comparable to an
energy meter, is intentionally not part of the current firmware. It was removed
from the practical project scope because it did not fit the concrete installation
concept: there was not enough suitable space for the controller close to the PV
generation side, and moving the controller farther away would have required long
measurement cable runs. Therefore, this feature was not goal-oriented for this
prototype, even though it remains technically possible in a different mechanical
and electrical setup.

The current default build receives the surplus value over MQTT from Home
Assistant. A pure local build is available by compiling with
\code{-D ENABLE_WIFI=0}; in that case the surplus value is read from the local
ADC input on GPIO34. There is also an optional compile-time fallback from MQTT to
the local ADC, but it is disabled by default because a floating ADC input would
produce unsafe meaningless values.

\newpage
\section{Result Against the Target Specification}

\begin{longtable}{p{0.30\textwidth}p{0.18\textwidth}p{0.42\textwidth}}
\toprule
Requirement & Status & Current implementation\\
\midrule
\endhead
Receive or measure PV surplus & \statuspartial &
Default build receives a signed feed-in value over MQTT. A local ADC path on
GPIO34 exists for offline builds or a future wired sensor. Integrated three-phase
measurement was described in the concept as a possible expansion, but it was not
implemented or tested because it was not practical for the actual project
installation.\\

Local switching decision & \statusdone &
\code{control.cpp} contains the state machine and decides the desired load state
on the ESP32. Network code only provides inputs and commands.\\

External switching output & \statusdone &
GPIO26 is the single firmware output for the load driver. It is initialized LOW
at boot and is only written through \code{hal::setLoad()}.\\

Modes OFF, MANUAL, AUTO & \statusdone &
The firmware implements all three modes. Priority is
\code{FAULT > OFF > MANUAL > AUTO}.\\

Tactile local user interface & \statusdone &
A physical button on GPIO25 is used. Short press cycles modes, long press toggles
the load in MANUAL or acknowledges faults.\\

Local status indication & \statuspartial &
Load, WiFi, fault and mode LEDs are implemented. The OLED interface exists as a
stub and is explicitly optional in firmware v1.\\

WiFi / Home Assistant integration & \statusdone &
WiFi, MQTT, retained state publishing and Home Assistant MQTT Discovery are
implemented.\\

Manual override & \statusdone &
In MANUAL mode, a long local button press or the Home Assistant manual switch can
request load ON/OFF. Fault and OFF still have higher priority.\\

Network-independent full operation & \statuspartial &
The local UI remains usable without WiFi. However, the default WiFi build needs a
fresh MQTT surplus value before any load can be energized, because a missing
surplus signal is treated as a safety fault. Full network-independent AUTO
operation would require the local ADC surplus source or a different local
measurement concept. The integrated three-phase measurement approach from the
original concept was excluded from this prototype for practical installation
reasons.\\

Final standalone three-phase measurement & \statusoutofscope &
The concept is documented, but no voltage/current measurement front end is
implemented in this project. In the concrete overall system concept, this
expansion was not
target-oriented for the prototype because the available installation space close
to the PV generation side was not sufficient. Placing the controller elsewhere
would have required long measurement cable runs, which would make the setup more
complex and less practical. Technically, integrated three-phase measurement would
still be possible, but it was not the most suitable goal for this project stage
and was therefore not tested.\\
\bottomrule
\end{longtable}

\newpage
\section{System Overview}

\subsection{Control chain}

The intended physical switching chain is:

\begin{lstlisting}
ESP32 GPIO26
  -> low-voltage transistor or MOSFET driver stage
  -> Ailao GSR2-1-10DA SSR input, 3-32 V DC
  -> SSR output switches 230 V AC contactor coil
  -> Heschen CT1-25 contactor switches the actual load
\end{lstlisting}

This chain is important because it separates the microcontroller from the
high-power circuit. The ESP32 only creates a 3.3 V logic signal. The actual load
current must be switched by a correctly rated contactor, not by the ESP32 and not
by a small GPIO-connected module.

\subsection{Main hardware components}

\begin{tabularx}{\textwidth}{p{0.26\textwidth}X}
\toprule
Component & Role in the system\\
\midrule
ESP32-32S / ESP32-WROOM-32 board & Main controller, local state machine, WiFi,
ADC readings and GPIO output.\\
Ailao GSR2-1-10DA SSR & Intended intermediate switching element. Its DC input is
controlled by the low-voltage driver stage; its AC output switches the contactor
coil.\\
Heschen CT1-25 contactor & Intended final switching element for the load circuit.
The concept uses the SSR only for the contactor coil, not for the final load.\\
10 kOhm NTC temperature sensor & Temperature monitoring input on GPIO35 through a
voltage divider.\\
Local button and LEDs & Tactile operation and local status feedback.\\
Optional OLED display & Not implemented in v1, but the firmware contains a
read-only interface for a later display implementation.\\
\bottomrule
\end{tabularx}

\subsection{Pin assignment}

\begin{longtable}{p{0.12\textwidth}p{0.22\textwidth}p{0.56\textwidth}}
\toprule
GPIO & Firmware name & Function\\
\midrule
\endhead
26 & \code{PIN_LOAD} & Digital output for the low-voltage SSR/contact driver.
Default LOW/OFF at boot.\\
25 & \code{PIN_BTN} & Local active-low multifunction button with internal
pull-up and software debounce.\\
32 & \code{PIN_MODE_A} & Reserved free GPIO, kept for possible future UI
extension.\\
33 & \code{PIN_MODE_B} & Reserved free GPIO.\\
34 & \code{PIN_PV} & ADC1 input for local PV surplus measurement in offline or
fallback builds.\\
35 & \code{PIN_NTC} & ADC1 input for the NTC voltage divider.\\
21 & \code{PIN_OLED_SDA} & I2C SDA for optional OLED.\\
22 & \code{PIN_OLED_SCL} & I2C SCL for optional OLED.\\
16 & \code{PIN_LED_LOAD} & Load status LED.\\
17 & \code{PIN_LED_WIFI} & WiFi/MQTT status LED.\\
18 & \code{PIN_LED_FAULT} & Fault LED.\\
27 & \code{PIN_LED_MODE} & Mode LED: off = OFF, blinking = AUTO, solid =
MANUAL.\\
\bottomrule
\end{longtable}

\subsection{Reason for the selected pin concept}

GPIO26 was chosen for the switching output because it is a regular output pin and
not one of the ESP32 flash pins or typical boot strapping pins. GPIO34 and GPIO35
are used for analog measurements because they are ADC1 input pins. This matters on
the ESP32 because ADC2 can conflict with WiFi operation. GPIO34 and GPIO35 are
input-only and have no internal pull-ups, which is acceptable for analog sensor
inputs but must be considered in the external wiring.

\newpage
\section{Firmware Architecture}

\subsection{Project structure}

\begin{lstlisting}
firmware/
  platformio.ini
  src/
    config.h
    main.cpp
    hal.h / hal.cpp
    temperature.h / temperature.cpp
    control.h / control.cpp
    net.h / net.cpp
    oled.h
    secrets.example.h
\end{lstlisting}

The firmware is built with PlatformIO for the \code{esp32dev} board and the
Arduino-ESP32 framework. The default external library is \code{PubSubClient} for
MQTT.

\subsection{Module responsibilities}

\begin{longtable}{p{0.24\textwidth}p{0.66\textwidth}}
\toprule
File or module & Responsibility\\
\midrule
\endhead
\code{config.h} & Central configuration: feature flags, pins, thresholds, ADC
constants, timing constants, fault policy and network tunables.\\
\code{main.cpp} & Composition root and non-blocking super-loop. It initializes the
system, gathers inputs, calls the control state machine and applies the resulting
load decision.\\
\code{hal.cpp} & Hardware abstraction layer. It owns direct GPIO and ADC access,
including the single writer for GPIO26.\\
\code{temperature.cpp} & Converts the NTC divider voltage into a temperature and
classifies NTC faults.\\
\code{control.cpp} & Pure decision logic. Implements modes, faults, latching,
hysteresis, delays and minimum dwell time.\\
\code{net.cpp} & WiFi, MQTT, Home Assistant Discovery, surplus input over MQTT and
remote commands. It never drives the load directly.\\
\code{oled.h} & Optional read-only display interface. In v1 it compiles to no-op
functions unless a later OLED implementation is added.\\
\bottomrule
\end{longtable}

\subsection{Reason for the modular split}

The split is a safety and verification decision. The code that receives MQTT
messages is intentionally separated from the code that can write GPIO26. A remote
command is only an input request. The state machine in \code{control.cpp} decides
whether the request is allowed, and \code{main.cpp} applies the decision through
\code{hal::setLoad()}. This makes it much easier to prove that faults, OFF mode
and boot behavior cannot be bypassed by network code.

\subsection{Main program flow}

The firmware uses a non-blocking super-loop. No steady-state \code{delay()} is
used. The loop repeatedly performs these tasks:

\begin{enumerate}
  \item feed the ESP32 task watchdog,
  \item service WiFi and MQTT,
  \item update PV and temperature samples,
  \item poll the local button and store short/long press events,
  \item every \code{CONTROL_TICK_MS}, process commands and evaluate the control
  state machine,
  \item write the desired load state through \code{hal::setLoad()},
  \item update LEDs and publish telemetry.
\end{enumerate}

This design is used because the controller must keep reacting to faults, button
presses and network events. Blocking delays would make the system less
predictable and could leave the load in an outdated state for too long.

\newpage
\section{Control Logic}

\subsection{Operating modes}

The firmware supports the three required operating modes:

\begin{itemize}
  \item \textbf{OFF}: the load is forced OFF.
  \item \textbf{MANUAL}: the manual latch decides whether the load should be ON
  or OFF.
  \item \textbf{AUTO}: the surplus value is compared against thresholds with
  hysteresis, delays and dwell time.
\end{itemize}

The selected mode is stored in \code{control::State::mode}. It is changed only by
explicit events: a local button short press or a Home Assistant mode command. The
selected mode is not overwritten by faults. This is intentional because Home
Assistant should continue to show a valid selected mode such as OFF, MANUAL or
AUTO while the fault itself is reported separately.

\subsection{Priority order}

Every control tick uses the same descending priority:

\begin{lstlisting}
FAULT -> OFF -> MANUAL -> AUTO
\end{lstlisting}

This means that a fault always switches the load off, even if the current selected
mode is MANUAL or AUTO. OFF also has priority over manual and automatic
operation. This simple priority rule is one of the main safety properties of the
firmware.

\subsection{AUTO mode}

AUTO mode is configured by default as follows:

\begin{itemize}
  \item switch-on threshold: \code{SURPLUS_ON_THRESHOLD_W = 1200 W},
  \item switch-off threshold: \code{SURPLUS_OFF_THRESHOLD_W = 800 W},
  \item switch-on delay: \code{SWITCH_ON_DELAY_MS = 10000 ms},
  \item switch-off delay: \code{SWITCH_OFF_DELAY_MS = 5000 ms},
  \item minimum dwell time before AUTO may switch on again:
  \code{MINIMUM_DWELL_TIME_MS = 30000 ms}.
\end{itemize}

The off threshold is lower than the on threshold. This hysteresis prevents rapid
switching around one single threshold. The delay timers require a continuous
condition: a short surplus spike is not enough to energize the contactor. The dwell
time reduces repeated cycling and mechanical stress. OFF transitions and faults
are not delayed, because safety must react immediately.

\subsection{Manual mode}

In MANUAL mode, the load follows a manual latch. Locally, the latch is toggled by
a long button press while the selected mode is MANUAL and no fault is active. Home
Assistant can also set the manual latch through its MQTT switch. MANUAL mode is
not dwell-gated; a deliberate local command should take effect immediately. A
fault or OFF mode still overrides the manual latch.

\subsection{Fault model}

Faults are represented as a bitmask so that several faults can be reported at the
same time.

\begin{longtable}{p{0.13\textwidth}p{0.25\textwidth}p{0.42\textwidth}p{0.12\textwidth}}
\toprule
Bit & Fault & Trigger & Clear behavior\\
\midrule
\endhead
\code{0x01} & PV invalid & PV input is not numeric, outside the plausible MQTT
range or outside the local ADC rail window. & Auto-clear on recovery.\\
\code{0x02} & NTC open & Divider node voltage is near 0 V, indicating a broken
or unplugged NTC path. & Healthy hold plus acknowledgement.\\
\code{0x04} & NTC short & Divider node voltage is near the supply rail,
indicating a short. & Healthy hold plus acknowledgement.\\
\code{0x08} & NTC range & Converted temperature is outside the configured
sensor range. & Healthy hold plus acknowledgement.\\
\code{0x10} & Overtemperature & NTC reads above \code{TEMPERATURE_LIMIT_C}. The
current firmware constant is 80 \textdegree C. & Auto-clear after cooling below
the clear threshold.\\
\code{0x20} & Mode selector & Reserved for an older selector-switch design. Not
raised in the current one-button UI. & Not used.\\
\code{0x40} & Sensor stale & PV or temperature sample path has stopped updating
within the configured freshness window. & Auto-clear on recovery.\\
\code{0x80} & Internal & A self-consistency condition is violated. & Manual
long-press override.\\
\bottomrule
\end{longtable}

The overtemperature limit is configurable in \code{config.h}. The current firmware
uses 80 \textdegree C with a 20 \textdegree C clear margin, so the overtemperature
fault can release after cooling below approximately 60 \textdegree C. This value
should be tuned after real thermal measurements of the driver stage and load
enclosure. The important point for the prototype is that overtemperature is
detected in the local control loop and immediately forces the load off.

\newpage
\section{Sensor Handling}

\subsection{PV surplus input}

The firmware supports two surplus input paths.

\textbf{MQTT path, default build:} Home Assistant publishes a signed feed-in power
value in watts. Positive values mean export or surplus. Negative values mean grid
import and are clamped to 0 W for the switching decision. Non-numeric values such
as \code{unknown} or \code{unavailable} are treated as invalid and force the load
off.

\textbf{Local ADC path:} GPIO34 maps a 0--3.3 V signal linearly to
0--3680 W. This is useful for a local test potentiometer or a future local sensor
front end. In a WiFi build, this path can be enabled as a fallback, but the default
is off because GPIO34 must not be left floating.

Both paths include freshness supervision. The MQTT surplus has a 90 s freshness
window because network sensors may update slowly. The local ADC path uses a much
shorter 500 ms freshness window.

\subsection{NTC temperature measurement}

The temperature sensor is a passive 10 kOhm NTC connected as a voltage divider:

\begin{lstlisting}
3V3 -- NTC -- GPIO35 -- 10 kOhm -- GND
\end{lstlisting}

With this topology, the node voltage rises as temperature rises. The firmware
uses:

\begin{lstlisting}
V_node = Vsupply * Rfixed / (Rntc + Rfixed)
Rntc   = Rfixed * (Vsupply / V_node - 1)
1/T    = 1/T0 + (1/Beta) * ln(Rntc / R0)
\end{lstlisting}

The Beta equation is sufficient for this prototype because the goal is not a
laboratory-grade thermometer, but a plausible thermal safety input. Open-circuit
and short-circuit checks are performed before the logarithmic conversion to avoid
invalid math and, more importantly, to ensure that a broken temperature sensor
cannot silently appear as a safe value.

\newpage
\section{Local User Interface}

\subsection{Button behavior}

The physical button on GPIO25 is active-low and uses the ESP32 internal pull-up.
It is debounced in software without blocking delays.

\begin{tabularx}{\textwidth}{p{0.30\textwidth}X}
\toprule
Action & Effect\\
\midrule
Short press & Cycle the selected mode: OFF, AUTO, MANUAL, OFF.\\
Long press in MANUAL & Toggle the manual load latch.\\
Long press while faulted & Acknowledge or override a fault clear where allowed.\\
\bottomrule
\end{tabularx}

Using a single multifunction button keeps the required tactile UI simple while
still covering all required actions. The trade-off is that the operator must learn
short press versus long press. This is acceptable for the prototype and avoids
using a larger physical control panel.

\subsection{LED behavior}

\begin{tabularx}{\textwidth}{p{0.22\textwidth}p{0.22\textwidth}X}
\toprule
LED & Pin & Meaning\\
\midrule
Load & GPIO16 & On when the commanded load state is ON.\\
WiFi & GPIO17 & Off = no WiFi, blinking = WiFi without MQTT, solid = MQTT
connected.\\
Fault & GPIO18 & Blinks while any fault is latched.\\
Mode & GPIO27 & Off = OFF, blinking = AUTO, solid = MANUAL.\\
\bottomrule
\end{tabularx}

\newpage
\section{Home Assistant and MQTT Integration}

\subsection{Configuration}

When \code{ENABLE_WIFI=1}, the firmware includes WiFi and MQTT support. Network
credentials and broker settings are not committed to the repository. The user
copies \code{firmware/src/secrets.example.h} to \code{firmware/src/secrets.h} and
fills in the SSID, password, broker host and MQTT topic.

The ESP32 subscribes to \code{MQTT_SURPLUS_TOPIC}, which defaults to
\code{pvctrl/surplus_w}. The expected payload is a signed number in watts.

\subsection{Home Assistant entities}

The firmware publishes Home Assistant MQTT Discovery configuration. A device named
\textit{PV Surplus Controller} appears with these entities:

\begin{longtable}{p{0.28\textwidth}p{0.18\textwidth}p{0.44\textwidth}}
\toprule
Entity & Direction & Purpose\\
\midrule
\endhead
Mode select & HA to ESP32 & Sets OFF, MANUAL or AUTO.\\
Manual Load switch & HA to ESP32 & Requests the manual latch state. Effective
only in MANUAL and only if no higher-priority fault exists.\\
Threshold ON number & HA to ESP32 & Sets the AUTO switch-on threshold.\\
Threshold OFF number & HA to ESP32 & Sets the AUTO switch-off threshold. Invalid
settings where OFF is not lower than ON are ignored.\\
Fault Reset button & HA to ESP32 & Requests fault acknowledgement.\\
Load binary sensor & ESP32 to HA & Reports the commanded load state.\\
Problem binary sensor & ESP32 to HA & Reports whether any fault is active.\\
Temperature sensor & ESP32 to HA & Reports the NTC temperature.\\
Surplus used sensor & ESP32 to HA & Reports the surplus value used by the
controller.\\
Fault Code sensor & ESP32 to HA & Reports the fault bitmask.\\
Uptime sensor & ESP32 to HA & Reports seconds since boot.\\
\bottomrule
\end{longtable}

\subsection{Reason for consume-once remote commands}

Mode and manual commands from Home Assistant are consumed once in
\code{main.cpp}. This prevents Home Assistant from overwriting a local button
press on every control tick. In a tick where both a remote and a local event
arrive, the firmware applies the remote command first and the local button second,
so the local action wins the tie. This supports the course requirement that the
device must remain genuinely operable at the physical device.

\newpage
\section{Safety Concept}

\subsection{Firmware safety measures}

\begin{itemize}
  \item \code{hal::initFailSafe()} is the first statement in \code{setup()} and
  drives GPIO26 LOW before any other initialization.
  \item GPIO26 is written only through \code{hal::setLoad()}.
  \item The controller is default-deny: the load is OFF unless MANUAL or AUTO
  explicitly request ON and no fault is active.
  \item The state machine checks faults before modes.
  \item Sensor freshness is supervised so that a frozen value cannot silently keep
  the load energized.
  \item The ESP32 task watchdog is armed. If the main loop stalls, the ESP32
  resets and the boot fail-safe drives the output LOW again.
  \item Remote commands never write GPIO directly.
\end{itemize}

\subsection{Hardware safety requirements}

The project is a proof of concept, not a certified mains-voltage product. The
following hardware points must be treated as mandatory for any real load test:

\begin{itemize}
  \item The ESP32 must never be connected to mains voltage.
  \item GPIO26 must drive only a low-voltage driver stage.
  \item The driver input needs an external pulldown resistor, for example
  10 kOhm, so the SSR input remains off during reset, brownout and boot.
  \item The SSR should switch only the 230 V contactor coil, not the final load.
  \item SSR off-state leakage must be checked to ensure the contactor drops out
  reliably.
  \item Work on 230 V or 400 V circuits requires qualified electrical work,
  correct fusing, enclosure, strain relief, clearances and thermal checks.
\end{itemize}

\subsection{Reason for using a contactor}

The contactor is used because high-power loads and especially three-phase loads
should not be switched by a microcontroller-side relay module. A contactor is
designed for load switching, has suitable contact ratings and provides a clearer
separation between the logic control side and the power circuit. The SSR is only
an intermediate element for the coil circuit.

\newpage
\section{Build and Operation}

\subsection{Build commands}

From the \code{firmware/} directory:

\begin{lstlisting}
pio run
pio run -t upload
pio device monitor
\end{lstlisting}

The default build enables WiFi and MQTT. It requires a local
\code{firmware/src/secrets.h}. For an offline build without network libraries:

\begin{lstlisting}
pio run -e esp32dev --build-flag "-D ENABLE_WIFI=0"
\end{lstlisting}

\subsection{Expected boot behavior}

On boot, GPIO26 is immediately forced LOW. After initialization, the firmware
starts the super-loop, updates sensors, evaluates the state machine every 100 ms
and periodically prints serial telemetry if serial debug output is enabled.

\newpage
\section{Verification Plan}

\subsection{Software-level checks}

Before connecting mains hardware, the following checks should be performed:

\begin{itemize}
  \item Compile the default or offline build with PlatformIO.
  \item Confirm by code search that GPIO26 is driven HIGH only inside
  \code{hal::setLoad()}.
  \item Confirm that \code{hal::initFailSafe()} is called before all other setup
  actions.
  \item Use the serial monitor to verify mode changes, fault codes, surplus value,
  temperature value and load decisions.
  \item Verify that no steady-state \code{delay()} calls exist in the control loop.
\end{itemize}

\subsection{Hardware verification sequence}

The test sequence follows the structure from the target specification.
The standalone local three-phase measurement from the original concept is not
part of this verification plan. It was a technically possible expansion idea, but
it was not implemented or tested because the real project installation did not
provide enough suitable space near the PV generation side and the resulting
measurement cable lengths would have made the solution impractical.

\begin{longtable}{p{0.14\textwidth}p{0.30\textwidth}p{0.46\textwidth}}
\toprule
Test & Purpose & Expected result\\
\midrule
\endhead
1.1 & GPIO26 control signal & GPIO26 produces a stable LOW/OFF and HIGH/ON logic
level during dry tests.\\
1.2 & Driver stage and SSR input & The low-voltage driver actuates the SSR input
reliably; the ESP32 remains on the low-voltage side.\\
1.3 & Single-phase AC load & A small test load follows the contactor command
without abnormal heating. Qualified mains setup required.\\
1.4 & Galvanic separation & No unsafe coupling is measured between ESP32 ground
and AC conductors. Qualified test required.\\
2.1 & Contactor coil actuation & The SSR output switches the CT1-25 coil and the
contactor pulls in and drops out reliably.\\
2.2 & Three-phase resistive load & All phases are switched by the contactor, not
by the ESP32 or SSR input side. Qualified test required.\\
2.3 & Repeated cycles & No contact welding, unintended retriggering or abnormal
temperature rise occurs.\\
2.4 & ESP32 power loss & Removing ESP32 power makes the contactor drop out and
the load de-energize.\\
3.1 & Surplus acquisition & MQTT or local ADC surplus is shown in serial output
and Home Assistant.\\
3.2 & AUTO switch-on threshold & In AUTO, load turns on only after surplus exceeds
the ON threshold continuously for the ON delay and dwell condition.\\
3.3 & Hysteresis switch-off & Load remains on until surplus falls below the lower
OFF threshold for the OFF delay.\\
3.4 & Manual override & In MANUAL, the local long press controls the manual latch
and takes priority over AUTO.\\
3.5 & Sensor failure & Invalid PV input, missing MQTT value or NTC fault forces
load OFF and reports a fault.\\
3.6 & Network independence & With local surplus input enabled, control continues
without WiFi. In the default MQTT build, WiFi or MQTT loss is intentionally a safe
OFF condition because the surplus signal is missing.\\
10.5 & Temperature monitoring & Heating or disconnecting the NTC triggers the
appropriate fault and switches the load OFF.\\
\bottomrule
\end{longtable}

\subsection{Suggested presentation demonstration}

For the exam presentation without slides, a compact demonstration can show:

\begin{enumerate}
  \item boot with GPIO26/load LED off,
  \item short button presses cycling OFF, AUTO and MANUAL,
  \item long press in MANUAL toggling the load output,
  \item simulated MQTT surplus causing AUTO switching after the delay,
  \item NTC open or short fault forcing the load off,
  \item Home Assistant showing mode, load state, surplus, temperature and fault
  code.
\end{enumerate}

This directly proves the parts that are central for the target specification:
local tactile control, local decision logic, safe output behavior and optional
network integration.

\newpage
\section{Known Limitations}

\begin{itemize}
  \item The current firmware is not a certified safety controller.
  \item Real mains-voltage and three-phase tests are hardware-dependent and must
  be performed by qualified persons.
  \item The OLED display is only a stub in firmware v1.
  \item The default WiFi build depends on an MQTT surplus value for load
  energization. Without a valid value, the firmware correctly chooses a fault and
  safe OFF state.
  \item Integrated three-phase measurement was considered in the concept, but it
  is not implemented or tested in this prototype because the required measurement
  placement and cable lengths were not practical for the real installation.
  \item There is no contactor auxiliary feedback input, so the firmware knows the
  commanded load state but not whether the contactor physically closed.
  \item ADC-based values should be calibrated per device before being used as
  precise measurements.
  \item Thresholds and mode commands are retained through Home Assistant/MQTT
  behavior, but there is no independent NVS configuration storage in the ESP32
  firmware.
\end{itemize}

\newpage
\section{Future Work}

The original concept described a possible final expansion stage in which the
controller becomes fully standalone through local three-phase measurement. This is
technically possible, but it was deliberately not pursued or tested in this
project because it did not fit the real installation concept. The available space
near the PV generation side was not sufficient for the additional measurement
hardware, and placing the controller in another location would have required long
measurement cable runs. For this prototype, using an external surplus value over
MQTT was therefore the more practical and goal-oriented solution.

If the mechanical and electrical installation conditions are changed in a future
version, such an expansion would require additional measurement hardware for all
three phases:

\begin{itemize}
  \item voltage measurement per phase,
  \item current measurement with current transformers,
  \item calculation of active, apparent and reactive power,
  \item detection of import/export direction,
  \item calculation of total surplus across all phases,
  \item optional battery or UPS supply for the low-voltage electronics.
\end{itemize}

On the firmware side, this can be integrated behind the existing surplus interface:
the control state machine already consumes a single surplus value in watts. That
means the measurement source can later change from MQTT or GPIO34 ADC to an
internal three-phase measurement module without redesigning the complete control
logic. This statement describes technical feasibility only; the measurement path
was not built or verified in the current project.

Other useful next steps are:

\begin{itemize}
  \item implement the OLED renderer,
  \item add a contactor feedback input,
  \item store a fault journal in non-volatile memory,
  \item add per-device ADC and NTC calibration,
  \item document measured thermal behavior of the SSR, driver stage and
  contactor enclosure.
\end{itemize}

\newpage
\section{Conclusion}

The implemented prototype satisfies the main goal of the course project: it shows
that an ESP32 can make a local switching decision for PV surplus load control,
while still offering Home Assistant integration for monitoring and configuration.
The firmware implements the required modes, tactile local operation, LED feedback,
fault handling, temperature monitoring, hysteresis and a fail-safe GPIO output.

The most important remaining work is the hardware verification of the mains
switching path. The final standalone three-phase measurement from the original
concept remains technically possible, but it was not a practical or
goal-oriented part of this project because of the available installation space and
the required cable lengths. The basic prototype is therefore best described as a
local-first ESP32 controller with an external surplus source and a carefully
separated low-voltage switching output.

\end{document}
